\chapter{十三个里程碑与实践方法}

\section{完整路线}

\begin{longtable}{p{0.11\textwidth}p{0.20\textwidth}p{0.60\textwidth}}
\toprule
版本 & 阶段 & 核心验收 \\
\midrule
v0.0 & 工程基线 & 可复现构建、测试分层、CI 与文档结构 \\
v0.1 & 复位与串口 & 自研 ROM 从复位向量输出串口协议 \\
v0.2 & 固件加载 & IDE PIO 加载自研 Stage 1，并覆盖错误状态 \\
v0.3 & Long Mode & 独立完成 16→32→64 位转换 \\
v0.4 & 内核加载 & 校验并加载 ELF64，交接 BootInfo \\
v0.5 & 内核基础 & GDT、IDT、TSS、异常与 panic \\
v0.6 & 内存管理 & 物理页、虚拟地址空间与内核堆 \\
v0.7 & 中断设备 & PIC、PIT、键盘与 IDE 驱动 \\
v0.8 & 用户边界 & Ring 3、用户 ELF 与系统调用 \\
v0.9 & 进程调度 & 抢占调度与进程生命周期 \\
v0.10 & 同步 IPC & 锁、等待队列、睡眠唤醒与管道 \\
v0.11 & 文件系统 & inode、目录、文件和持久化 \\
v1.0 & 用户环境 & Shell、基础命令和整机回归 \\
\bottomrule
\end{longtable}

\section{每个阶段的固定问题}

\begin{enumerate}[label=\arabic*.]
  \item 硬件或上阶段提供什么输入状态？
  \item 本阶段新建立哪些不变量？
  \item 输出交给谁，地址和 ABI 是什么？
  \item 正常路径怎样观察？
  \item 超时、损坏和非法输入怎样失败？
  \item 哪些单元、集成、随机和整机测试保存证据？
\end{enumerate}

\section{代码走读方法}

先从链接脚本和入口找最终地址，再沿调用链跟踪寄存器和内存所有权；遇到端口
访问就回到设备寄存器协议；遇到跨层调用就写下输入、输出与失败状态；最后对照
测试，确认每个关键分支都能被外部证据区分。

\section{学习记录}

每完成一个机制，写一页实验记录：假设、最小实现、观察结果、失败案例、根因和
后续限制。不要只保存成功截图。真正提高定位能力的是“为什么没有输出”“哪一
个字节错了”“哪个状态位没有到达”。
